\section{Problem and Security Model}
\label{sec:threat-model}

\paragraph{System and enforcement boundary.}
A user gives an agent an authenticated task $q$. The agent may read messages,
documents, webpages, or database records before proposing a structured tool
call $c=(t,\bar{x})$. The large language model (LLM) is an untrusted control-plane component: its plan,
effect description, and arguments are proposals rather than authority.
Executing the totalized call from pre-state $s$ may commit several externally
visible effects, such as creating an event, inviting principals, or changing
visibility. Enforcement occurs after argument totalization and immediately
before commit. The executor may commit only the exact call associated with the
monitor's decision.

\paragraph{Threat model.}
The adversary controls content returned by one or more untrusted reads. It may
embed instructions or attacker-chosen values, induce the agent to add a
resource, target, payload, operation, or effect, and carry this influence across
a multi-step trajectory; recent studies show that agent skills leak
credentials through exactly this channel~\cite{chen2026credentials}. We
assume the agent may follow the injected
instruction completely; security does not rely on model alignment. The
adversary cannot modify the authenticated task, forge trusted provenance,
change a registered descriptor or policy, compromise the monitor, or replace a
checked call before execution. It also cannot replace the registered tool
implementation without triggering revalidation.

\paragraph{Representation and commit objective.}
Let $u=(t,\bar{x},s,h)$ denote an execution context containing a tool,
totalized arguments, trusted pre-state, and provenance evidence. Let $\mu(u)$
denote its committed effects. A policy context $P$ may additionally depend on
trusted execution metadata $\gamma(u)$, including effect provenance. Writing
$\omega(u)=(\mu(u),\gamma(u))$, the policy context $P$ induces an ideal
per-commit predicate $I_P(\omega(u))$. The desired commit condition is $I_P(\omega(u))=1$.

The explicit-authority experiment instantiates $P$ with authenticated subject
identity, resource ACLs, held capabilities, and attenuated delegations. The
earlier source-relation diagnostic uses effect-set containment only as a finite
interface-discrimination stress relation. In both settings, the descriptor is
an observation contract rather than a source of authority: the authenticated
user, harness, or external policy component supplies $P$, and the planner
cannot enlarge it. A separate AgentDojo integration implements a narrower
provenance-origin predicate. It tests mediation and auditability, not
organization-wide entitlement.

\paragraph{Registration-time trust assumptions.}
Offline validation assumes that the copied sandbox faithfully runs the
versioned implementation, the source execution oracle observes the exercised
committed transition, interventions are reproducible, and the
policy-separation relation is supplied independently of the candidate
descriptor. These privileged observations generate registration evidence and
do not sit on the deployed commit path. Failed or unresolved executions remain
explicit and cannot justify dropping a field.

\paragraph{Runtime enforcement TCB.}
Runtime trusts the descriptor--implementation binding, authority and policy
state, the provenance adapter where used, the monitor, and the exact-call
executor. State-dependent instantiation additionally requires a fresh trusted
pre-commit snapshot. The conditional guarantee assumes complete mediation and
check--use consistency for both the call and policy snapshot. If the planner
can alter authority, relabel provenance, replace the checked call, or race the
trusted state, it has crossed the stated boundary. Section~\ref{sec:limitations}
discusses semantic drift, concurrency, remote effects, and output-only goals.
